\documentclass[a4paper,12pt]{article}

\usepackage[T2A]{fontenc}
\usepackage[utf8]{inputenc}
\usepackage[russian]{babel}
\usepackage{amsmath,amssymb,mathtools}
\usepackage{PTSans}
\renewcommand{\familydefault}{\sfdefault}
\usepackage{icomma}
\usepackage[a4paper,margin=20mm]{geometry}
\usepackage{microtype}
\usepackage{tikz}
\usetikzlibrary{arrows.meta,positioning,shapes.geometric}
\usepackage{tabularx,booktabs}
\usepackage{enumitem}
\usepackage{fancyhdr}
\usepackage{setspace}
\usepackage{xcolor}
\usepackage[hidelinks]{hyperref}

\definecolor{accent}{HTML}{2F6F73}
\definecolor{darkaccent}{HTML}{204B4E}
\definecolor{textgray}{HTML}{303638}
\definecolor{softgray}{HTML}{6D7779}
\definecolor{lightline}{HTML}{D6DFDF}

\geometry{headheight=25pt}
\onehalfspacing
\setlength{\parindent}{0pt}
\setlength{\parskip}{0.45em}
\setlist[itemize]{leftmargin=1.55em,itemsep=0.25em,topsep=0.3em}
\setlist[enumerate]{leftmargin=1.7em,itemsep=0.5em,topsep=0.35em}

\pagestyle{fancy}
\fancyhf{}
\fancyhead[L]{\small\color{softgray}Чек-лист по математике}
\fancyhead[R]{\small\color{softgray}29.09.26}
\fancyfoot[C]{\small\color{softgray}\thepage}
\renewcommand{\headrulewidth}{0.3pt}
\renewcommand{\headrule}{\hbox to\headwidth{\color{lightline}\leaders\hrule height \headrulewidth\hfill}}

\newcommand{\topic}[1]{%
  \vspace{0.55em}%
  {\large\bfseries\color{darkaccent}#1}\par
  \vspace{0.15em}\color{lightline}\rule{\linewidth}{0.6pt}\color{textgray}\vspace{0.25em}
}
\newcommand{\key}[1]{\textbf{\color{darkaccent}#1}}
\newcommand{\checkitem}{\(\square\)\;}
\newcommand{\nodop}{\mathop{\text{НОД}}\nolimits}
\newcommand{\nokop}{\mathop{\text{НОК}}\nolimits}

\begin{document}
\color{textgray}

% -------------------------------------------------
% TITLE PAGE
% -------------------------------------------------
\begin{titlepage}
\thispagestyle{empty}
\begin{center}
\vspace*{2.8cm}

{\Large\color{softgray}Математика}\par
\vspace{0.9cm}
{\fontsize{30}{36}\selectfont\bfseries\color{darkaccent}Чек-лист}\par
\vspace{0.35cm}
{\LARGE\bfseries Обыкновенные дроби}\par
\vspace{0.2cm}
{\large сокращение, общий знаменатель и десятичная запись}\par

\vspace{1.5cm}
\color{lightline}\rule{0.72\linewidth}{0.7pt}\par
\vspace{1.1cm}

{\large Ксения Васильченко}\par
\vspace{0.35cm}
{\normalsize 29.09.26}\par

\vfill
{\small\color{softgray}Лёвин Артём Александрович эксклюзивно для Ксении Васильченко}\par
\vspace{0.8cm}

\begin{tikzpicture}[x=1cm,y=1cm]
  \draw[accent!60,line width=1.25pt]
    (-6.2,0) .. controls (-3.6,1.0) and (-1.7,-0.9) .. (0,0)
    .. controls (1.8,0.95) and (3.7,-0.6) .. (6.2,0.2);
  \fill[accent!35] (-4.0,0.52) circle (1.7pt);
  \fill[accent!35] (3.6,-0.20) circle (1.7pt);
\end{tikzpicture}
\end{center}
\end{titlepage}

% -------------------------------------------------
% CHECKLIST
% -------------------------------------------------
\section*{Главное после занятия}

\topic{1. Вычислительный фундамент}
\begin{itemize}
  \item \checkitem Умею складывать и вычитать десятичные дроби, записывая запятую под запятой.
  \item \checkitem Умею умножать десятичные дроби столбиком: сначала работаю как с целыми, затем отделяю справа столько знаков, сколько их было после запятой у обоих множителей вместе.
  \item \checkitem Умею делить десятичную дробь на натуральное число и правильно ставить запятую в частном.
  \item \checkitem Помню действия с обыкновенными дробями: при умножении перемножаю числители и знаменатели; при делении умножаю на обратную дробь.
  \item \checkitem Перевожу смешанную дробь в неправильную и обратно.
\end{itemize}

\topic{2. Сокращение дробей}
Дробь можно сокращать только тогда, когда \key{числитель и знаменатель делятся на один и тот же ненулевой множитель}.

\[
\frac{14}{28}=\frac{2\cdot7}{2\cdot2\cdot7}=\frac12,
\qquad
\frac{10}{35}=\frac{2\cdot5}{5\cdot7}=\frac27.
\]

\key{Универсальный способ:} разложить числитель и знаменатель на простые множители, затем сократить общие множители.

\key{Быстрый способ:} если общий делитель виден сразу, можно разделить числитель и знаменатель непосредственно на него.

\topic{3. Приведение дроби к заданному знаменателю}
Пусть дано \(\dfrac{a}{b}\), а нужен знаменатель \(D\). Это возможно, если \(D\) делится на \(b\) без остатка.

\[
k=\frac{D}{b},\qquad
\frac{a}{b}=\frac{a\cdot k}{b\cdot k}.
\]

Например,
\[
\frac12=\frac{1\cdot28}{2\cdot28}=\frac{28}{56},
\qquad
\frac27=\frac{2\cdot8}{7\cdot8}=\frac{16}{56}.
\]

\key{Важно:} числитель и знаменатель всегда домножаются на одно и то же число. Менять только знаменатель нельзя.

\topic{4. Можно ли получить заданный знаменатель?}
Чтобы проверить, можно ли привести дробь \(\dfrac{a}{b}\) к знаменателю \(D\), вычислите \(D:b\).

\begin{itemize}
  \item Если получилось целое число, приведение возможно.
  \item Если деление идёт с остатком, получить такой знаменатель домножением нельзя.
\end{itemize}

\[
48:3=16 \Rightarrow \frac13=\frac{16}{48},
\qquad
48:9\notin\mathbb Z \Rightarrow \frac49\text{ к знаменателю }48\text{ не приводится}.
\]

\topic{5. Общий и наименьший общий знаменатель}
Для нескольких дробей общий знаменатель должен делиться на каждый исходный знаменатель. Самый удобный выбор --- \key{НОК знаменателей}.

Например, для \(\dfrac{5}{12}\) и \(\dfrac{7}{18}\):
\[
\nokop(12,18)=36,
\qquad
\frac{5}{12}=\frac{15}{36},
\qquad
\frac{7}{18}=\frac{14}{36}.
\]
После этого дроби можно складывать или вычитать:
\[
\frac{5}{12}-\frac{7}{18}=\frac{15}{36}-\frac{14}{36}=\frac1{36}.
\]

\topic{6. НОД и НОК: короткое напоминание}
\begin{itemize}
  \item \key{НОД} --- произведение общих простых множителей в наименьших степенях.
  \item \key{НОК} --- произведение всех нужных простых множителей в наибольших степенях.
\end{itemize}

Для \(12=2^2\cdot3\) и \(15=3\cdot5\):
\[
\nodop(12,15)=3,
\qquad
\nokop(12,15)=2^2\cdot3\cdot5=60.
\]

\topic{7. Обыкновенная дробь как конечная десятичная}
Сначала \key{сократите дробь}. После сокращения конечная десятичная запись существует тогда и только тогда, когда в разложении знаменателя есть только простые множители \(2\) и \(5\):
\[
b=2^m5^n.
\]

Это означает, что знаменатель можно домножить до \(10\), \(100\), \(1000\), \(10000\) и т.~д.

\[
\frac{3}{5}=\frac{6}{10}=0{,}6,
\qquad
\frac{12}{25}=\frac{48}{100}=0{,}48,
\qquad
\frac{7}{20}=\frac{35}{100}=0{,}35.
\]

А дроби \(\dfrac13\) и \(\dfrac7{12}\) конечной десятичной записи не имеют, потому что после сокращения в знаменателе остаётся множитель \(3\).

\topic{8. Знаменатель уже равен $10^n$}
Если знаменатель равен \(10\), \(100\), \(1000\), \(10000\), то запятая в числителе сдвигается влево на число нулей в знаменателе:
\[
\frac{123}{10}=12{,}3,\qquad
\frac{123}{100}=1{,}23,\qquad
\frac{123}{1000}=0{,}123,\qquad
\frac{123}{10000}=0{,}0123.
\]

\topic{9. Самопроверка перед ответом}
\begin{itemize}
  \item \checkitem Дробь сокращена, если это возможно.
  \item \checkitem При домножении изменены и числитель, и знаменатель.
  \item \checkitem Новый знаменатель действительно получен из старого умножением на целое число.
  \item \checkitem При сложении/вычитании знаменатели сначала стали одинаковыми.
  \item \checkitem Для конечной десятичной дроби критерий применён после сокращения.
  \item \checkitem В десятичной записи поставлено нужное число нулей перед значащими цифрами.
\end{itemize}

\topic{10. Типичные ошибки}
\begin{itemize}
  \item Сокращать слагаемые вместо множителей, например пытаться сокращать в \(\dfrac{a+b}{b}\).
  \item Домножать только знаменатель и тем самым менять значение дроби.
  \item Искать общий знаменатель, который не делится на один из знаменателей.
  \item Считать любую дробь со знаменателем, оканчивающимся на ноль, десятичной: например, \(\dfrac7{20}\) сначала нужно привести к знаменателю \(100\).
  \item Проверять множители знаменателя до сокращения дроби. Например, \(\dfrac36=\dfrac12=0{,}5\), хотя в числе \(6\) есть множитель \(3\).
\end{itemize}

% -------------------------------------------------
% FLOWCHART PAGE
% -------------------------------------------------
\clearpage
\thispagestyle{fancy}
\begin{center}
{\LARGE\bfseries\color{darkaccent}Алгоритм: привести дроби к общему знаменателю}
\end{center}
\vspace{1cm}

\begin{center}
\begin{tikzpicture}[
  node distance=10mm and 12mm,
  startstop/.style={ellipse,draw=accent,line width=0.9pt,minimum width=46mm,minimum height=11mm,align=center,fill=accent!4},
  process/.style={rectangle,rounded corners=2mm,draw=accent,line width=0.9pt,text width=105mm,minimum height=12mm,align=center,fill=accent!3},
  decision/.style={diamond,aspect=2.0,draw=accent,line width=0.9pt,text width=32mm,align=center,inner sep=2pt,fill=accent!4},
  arrow/.style={-{Latex[length=3mm]},line width=0.9pt,color=darkaccent}
]
  \node[startstop] (start) {Даны две или несколько дробей};
  \node[process,below=of start] (reduce) {Сократите каждую дробь, если это возможно};
  \node[process,below=of reduce] (lcm) {Найдите \(D=\nokop\) знаменателей};
  \node[process,below=of lcm] (factor) {Для каждой дроби \(\dfrac{a}{b}\) найдите дополнительный множитель \(k=D:b\)};
  \node[process,below=of factor] (multiply) {Домножьте и числитель, и знаменатель: \(\dfrac{a}{b}=\dfrac{a\cdot k}{b\cdot k}\)};
  \node[decision,below=12mm of multiply] (check) {Все знаменатели равны \(D\)?};
  \node[startstop,below=13mm of check] (finish) {Готово: выполняйте нужное действие с дробями};
  \node[process,right=9mm of check,text width=38mm] (fix) {Проверьте \(D\), деление \(D:b\) и домножение};

  \draw[arrow] (start) -- (reduce);
  \draw[arrow] (reduce) -- (lcm);
  \draw[arrow] (lcm) -- (factor);
  \draw[arrow] (factor) -- (multiply);
  \draw[arrow] (multiply) -- (check);
  \draw[arrow] (check) -- node[right,font=\small]{да} (finish);
  \draw[arrow] (check.east) -- node[above,font=\small]{нет} (fix.west);
  \draw[arrow] (fix.north) |- (factor.east);
\end{tikzpicture}
\end{center}
\clearpage

% -------------------------------------------------
% TRAINING
% -------------------------------------------------
\section*{Тренировочный блок --- Путь А}
\textit{10 заданий. Решения оформляйте отдельно; здесь оставлены только условия.}

\begin{enumerate}[label=\textbf{\arabic*.}]
  \item Вычислите:
  \[
  4{,}7-2{,}38+0{,}65.
  \]

  \item Вычислите и выделите целую часть, если получится неправильная дробь:
  \[
  \frac34+\frac56.
  \]

  \item Переведите смешанную дробь в неправильную и вычислите:
  \[
  2\frac13\cdot\frac9{14}.
  \]

  \item Найдите \(\nodop(18,24)\) и \(\nokop(18,24)\).

  \item Сократите дробь \(\dfrac{42}{63}\), затем приведите полученную дробь к знаменателю \(18\).

  \item Определите, можно ли дробь \(\dfrac7{15}\) привести:
  \begin{enumerate}[label=\alph*)]
    \item к знаменателю \(90\); если можно, приведите;
    \item к знаменателю \(80\).
  \end{enumerate}

  \item Приведите дроби \(\dfrac5{12}\) и \(\dfrac7{18}\) к наименьшему общему знаменателю.

  \item Вычислите:
  \[
  \frac5{12}-\frac7{18}.
  \]

  \item Представьте в виде конечных десятичных дробей:
  \[
  \frac7{40},\qquad \frac{13}{125}.
  \]

  \item Какие из дробей имеют конечную десятичную запись? Для тех, которые имеют, запишите её:
  \[
  \frac7{24},\qquad \frac9{40},\qquad \frac{13}{50},\qquad \frac5{18}.
  \]
\end{enumerate}

% -------------------------------------------------
% ANSWERS
% -------------------------------------------------
\clearpage
\section*{Ответы}
\renewcommand{\arraystretch}{1.45}
\begin{center}
\begin{tabularx}{\linewidth}{@{}>{\bfseries}p{12mm}X@{}}
\toprule
№ & Ответ \\
\midrule
1 & \(2{,}97\). \\
2 & \(\dfrac{19}{12}=1\dfrac7{12}\). \\
3 & \(\dfrac32=1\dfrac12\). \\
4 & \(\nodop(18,24)=6\), \(\nokop(18,24)=72\). \\
5 & \(\dfrac{42}{63}=\dfrac23=\dfrac{12}{18}\). \\
6 & а) можно: \(\dfrac7{15}=\dfrac{42}{90}\); б) нельзя, так как \(80:15\) не является целым числом. \\
7 & \(\nokop(12,18)=36\); \(\dfrac5{12}=\dfrac{15}{36}\), \(\dfrac7{18}=\dfrac{14}{36}\). \\
8 & \(\dfrac1{36}\). \\
9 & \(\dfrac7{40}=0{,}175\), \(\dfrac{13}{125}=0{,}104\). \\
10 & Конечные: \(\dfrac9{40}=0{,}225\), \(\dfrac{13}{50}=0{,}26\). Дроби \(\dfrac7{24}\) и \(\dfrac5{18}\) конечной десятичной записи не имеют. \\
\bottomrule
\end{tabularx}
\end{center}

\vfill
{\small\color{softgray}Проверка: после сокращения конечная десятичная запись возможна только при знаменателе вида \(2^m5^n\).}

\end{document}
